\definecolor{ISEL}{RGB}{163, 38, 11}

\begin{titlepage}  
    \begin{center}
    
        \vspace*{\fill}   
        
        \huge \textbf{Circuitos Aritméticos e Lógicos}
        
        \large Relatório do 2º Trabalho
        \vspace{1.25in}
        
        \normalsize
        
        \textbf{Desenvolvimento de uma Unidade Aritmética e Lógica em VHDL}
        \vspace{5mm}

        Este trabalho apresenta o desenvolvimento de uma ALU (\textit{Arithmetic Logic Unit}) em VHDL, com foco na implementação de operações aritméticas e lógicas básicas em 4 bits. A ALU foi projetada e testada na placa DE10-Lite da \textit{Intel}.
        \vspace{5mm}

        Através deste projeto, procura-se não apenas entender o funcionamento deste tipo de unidade, mas também aplicar teorias da computação e lógica digital, num ambiente prático de desenvolvimento...
        \vspace{1in}
        
        Trabalho realizado por: 
        
        \textbf{Gustavo Costa | Nº 52808}
        
        \textbf{Rafael Pereira | Nº 52880}
        \vspace{1in}
        
        Turma: \textbf{LEIC15D}
        
        Docente: \textbf{Prof.º Mário Véstias}
        \vspace{1.25in}
        
        Licenciatura em Engenharia Informática e de Computadores
        
        Lógica e Sistemas Digitais (LSD)
        
        2024 / 2025 - Semestre de Inverno
        
        \vspace*{\fill}%
        
    \end{center}
\end{titlepage}